\chapter{Referencia rápida de la API de \pyhydra}
\label{app:api}

Este anexo es una referencia rápida de las funciones y clases principales de \pyhydra organizadas por módulo. Para cada componente se indica el propósito, los parámetros principales con sus tipos y valores por defecto, y el tipo de retorno. La relación se ha revisado contra el árbol de código de \pyhydra incluido con la tesis; por tanto, describe funciones y clases presentes en el paquete, no una interfaz idealizada ni funcionalidades planificadas. La documentación completa se encuentra en el capítulo de manual de uso (\cref{ch:manual}) y en los notebooks correspondientes (\cref{app:notebooks}).

\newcommand{\apisig}[1]{%
  \begingroup
  \setlength{\fboxsep}{3pt}%
  \colorbox{codebg}{\parbox{0.96\linewidth}{\raggedright\ttfamily\footnotesize #1}}%
  \endgroup
}

% ============================================================
\section{\texttt{pyhydra.data\_sources}}
% ============================================================

\subsection{Precipitación (\texttt{rainfall})}

\begin{description}[style=nextline,font=\ttfamily]
  \item[\apisig{download\_era5(folder\_out, api\_key, url, area, variables\_list, years,\\ months=range(1,13), max\_workers=5,\\ file\_format=\textquotesingle{}netcdf\textquotesingle{}, frequency=\textquotesingle{}hourly\textquotesingle{})}]
  Descarga ERA5 mediante la API CDS en paralelo por año y mes. \texttt{area} sigue el formato \texttt{[N, W, S, E]}. \texttt{frequency} acepta \texttt{\textquotesingle{}hourly\textquotesingle{}} o \texttt{\textquotesingle{}monthly\textquotesingle{}}. Sin valor de retorno (guarda ficheros NetCDF en \texttt{folder\_out}).

  \item[\apisig{GPMDownloader()}]
  Clase para descarga y apertura de productos GPM-IMERG vía \texttt{earthaccess} (autentica contra NASA Earthdata en el constructor). Métodos: \path{.set_region(points=None, lat_bounds=None, lon_bounds=None)}, \path{.search(start_date, end_date, resolution='hourly')} y \path{.open_dataset(results, output_folder=...)}. Devuelve \texttt{xarray.Dataset}.

  \item[\apisig{PERSSIANDownloader(lon=None, lat=None, lon\_min=None, lon\_max=None,\\ lat\_min=None, lat\_max=None, path\_output=\textquotesingle{}output\textquotesingle{}, max\_workers=2)}]
  Clase para descarga de PERSIANN-CCS diario desde el servidor FTP de UCI, en modo punto (\texttt{lon}/\texttt{lat}) o bounding-box (\texttt{lon\_min}/\texttt{lon\_max}/\texttt{lat\_min}/\texttt{lat\_max}). Método principal: \texttt{.download\_daily(start\_date, end\_date)}. Sin valor de retorno.

  \item[\apisig{download\_aemet\_daily\_data(path\_output, api\_key, start\_date\_str,\\ end\_date\_str, interval\_days=15)}]
  Descarga observaciones diarias de todas las estaciones AEMET OpenData. \texttt{interval\_days} controla el tamaño del bloque de descarga (límite de la API). Sin valor de retorno.

  \item[\apisig{AEMETDownloader(netcdf\_folder, api\_key=None, stations\_shapefile=None)}]
  Función (no clase) que construye un widget Jupyter interactivo para selección y descarga de estaciones AEMET; en modo directo requiere \texttt{api\_key}. Devuelve el widget \texttt{ipywidgets}.

  \item[\apisig{AemetCSVLoader(folder\_path)}]
  Carga y filtra de calidad series descargadas por \texttt{AEMETDownloader}. Métodos: \texttt{.load\_station\_data()} (lee \texttt{selected\_stations.csv}, devuelve \texttt{pd.DataFrame}), \texttt{.load\_series\_data(variable)} (concatena las series de la variable indicada) y \texttt{.analyze\_series\_quality()}.

  \item[\apisig{download\_synop(station\_id, start\_date, end\_date, progress=None,\\ cancel\_flag=None)}]
  Descarga datos SYNOP de OGIMET para una estación y un período. Devuelve \texttt{pd.DataFrame} con variables meteorológicas.

  \item[\apisig{OGIMETDownloader(stations\_csv=None, zoom=5, center=(40.0,-3.5),\\ max\_markers=None, output\_folder=None)}]
  Función (no clase) que construye un widget interactivo con mapa Leaflet para descarga por lotes desde OGIMET.

  \item[\apisig{OgimetCSVLoader(folder\_path=None, create=True)}]
  Carga y consolida series descargadas por \texttt{OGIMETDownloader}. Métodos principales:
  \texttt{.load\_station\_data()}, \texttt{.load\_series\_data(variable)}
  y \texttt{.analyze\_series\_quality()}. Sigue el mismo patrón que \texttt{AemetCSVLoader}.

  \item[\apisig{download\_meteostat\_series(station\_id, start\_date, end\_date,\\ frequency=\textquotesingle{}daily\textquotesingle{}, normalize=True, station\_name=None)}]
  Descarga series de una estación Meteostat por su código identificador. \texttt{frequency} acepta \texttt{\textquotesingle{}daily\textquotesingle{}} u \texttt{\textquotesingle{}hourly\textquotesingle{}}. Devuelve \texttt{pd.DataFrame} con columnas de temperatura, precipitación, humedad, viento y presión.

  \item[\apisig{download\_best\_meteostat\_series(lat, lon, start\_date, end\_date,\\ radius=150\_000, frequency=\textquotesingle{}daily\textquotesingle{})}]
  Combina búsqueda, puntuación y descarga en una sola llamada. \texttt{radius} está en metros. Selecciona la estación con mejor cobertura de datos dentro del radio. Devuelve tupla \texttt{(series, station\_id, ranking\_df)}.

  \item[\apisig{get\_meteostat\_stations\_nearby(lat, lon, radius=150\_000, limit=100)}]
  Lista estaciones Meteostat próximas a un punto. \texttt{radius} en metros. Devuelve \texttt{pd.DataFrame} con metadatos de estaciones.

  \item[\apisig{score\_meteostat\_stations(stations\_df, start\_date, end\_date,\\ frequency=\textquotesingle{}daily\textquotesingle{}, variables=None, distance\_weight=0.10)}]
  Puntúa estaciones según fracción de datos disponibles en el período (ponderada con la distancia). Devuelve \texttt{pd.DataFrame} con columna \texttt{score}.

  \item[\apisig{find\_station\_by\_wmo(wmo\_code)}]
  Busca una estación Meteostat por su código WMO. Devuelve \texttt{pd.DataFrame}.

  \item[\apisig{MeteostatDownloader(zoom=5, center=(40.0,-3.5), radius=250\_000,\\ max\_markers=300, output\_folder=None)}]
  Función (no clase) que construye un widget JupyterLab con mapa Leaflet para exploración y descarga interactiva de series Meteostat; \texttt{radius} en metros.
\end{description}

\subsection{Caudal (\texttt{river\_discharge})}

\begin{description}[style=nextline,font=\ttfamily]
  \item[\apisig{download\_glofas(area, years, output\_dir,\\ dataset=\textquotesingle{}cems-glofas-historical\textquotesingle{}, system\_version=\textquotesingle{}version\_4\_0\textquotesingle{},\\ hydrological\_model=\textquotesingle{}lisflood\textquotesingle{}, product\_type=\textquotesingle{}consolidated\textquotesingle{},\\ months=None, variable=\textquotesingle{}river\_discharge\_in\_the\_last\_24\_hours\textquotesingle{},\\ file\_format=\textquotesingle{}netcdf4\textquotesingle{}, api\_key=None)}]
  Descarga reanálisis de caudal GloFAS desde el CDS de Copernicus. \texttt{area} en formato \texttt{[N, W, S, E]}. \texttt{dataset} puede ser \texttt{\textquotesingle{}cems-glofas-historical\textquotesingle{}}, \texttt{\textquotesingle{}cems-glofas-reforecast\textquotesingle{}} o \texttt{\textquotesingle{}cems-glofas-forecast\textquotesingle{}}. Devuelve lista de rutas a ficheros descargados.

  \item[\apisig{read\_glofas\_nc(filepath, lat, lon, lat\_col=\textquotesingle{}latitude\textquotesingle{},\\ lon\_col=\textquotesingle{}longitude\textquotesingle{})}]
  Lee un fichero NetCDF de GloFAS y extrae la serie temporal en el punto más próximo a \texttt{(lat, lon)}. Devuelve \texttt{pd.Series}.

  \item[\apisig{read\_grdc(filepath)}]
  Lee un fichero de texto GRDC en formato estándar. Devuelve \texttt{pd.DataFrame} con índice de fechas.

  \item[\apisig{read\_grdc\_metadata(filepath)}]
  Extrae los metadatos del encabezado de un fichero GRDC. Devuelve \texttt{dict}.

  \item[\apisig{read\_grdc\_folder(folder, pattern=\textquotesingle{}*.day\textquotesingle{})}]
  Lee todos los ficheros GRDC de un directorio. Devuelve \texttt{dict} mapeando nombre de estación a \texttt{pd.DataFrame}.

  \item[\apisig{analyze\_grdc\_quality(df)}]
  Calcula estadísticos de calidad de una serie GRDC: fracción de datos válidos, saltos, valores fuera de rango. Devuelve \texttt{dict}.

  \item[\apisig{download\_usgs(site\_no, start\_date, end\_date,\\ parameter\_cd=\textquotesingle{}00060\textquotesingle{}, units=\textquotesingle{}metric\textquotesingle{}, max\_retries=3)}]
  Descarga datos de aforos USGS mediante la API NWIS. \texttt{site\_no} puede ser cadena o lista. \texttt{parameter\_cd=\textquotesingle{}00060\textquotesingle{}} corresponde a caudal. \texttt{units=\textquotesingle{}metric\textquotesingle{}} devuelve m$^3$/s (por defecto); \texttt{\textquotesingle{}imperial\textquotesingle{}} devuelve ft$^3$/s. Devuelve \texttt{pd.DataFrame}.

  \item[\apisig{get\_usgs\_site\_info(site\_no)}]
  Obtiene metadatos de una o varias estaciones USGS. Devuelve \texttt{pd.DataFrame}.

  \item[\apisig{search\_usgs\_sites(bbox, parameter\_cd=\textquotesingle{}00060\textquotesingle{})}]
  Busca estaciones USGS dentro de un bounding box \texttt{(west, south, east, north)} en grados decimales. Devuelve \texttt{pd.DataFrame}.
\end{description}

\subsection{Cambio climático y suelos}

\begin{description}[style=nextline,font=\ttfamily]
  \item[\apisig{download\_CDS\_CMIP6(url, api\_key, start\_date, end\_date,\\ temporal\_resolution, model, experiments, variables,\\ download\_base\_dir, area, combinations\_csv=None,\\ max\_workers=5, max\_retries=3)}]
  Descarga proyecciones CMIP6 del CDS para una combinación de modelo, experimento, variable, período y dominio. \texttt{experiments} puede incluir \texttt{\textquotesingle{}historical\textquotesingle{}}, \texttt{\textquotesingle{}ssp245\textquotesingle{}}, \texttt{\textquotesingle{}ssp585\textquotesingle{}}, etc. Sin valor de retorno (guarda ficheros en \texttt{download\_base\_dir}).

  \item[\apisig{get\_dataset\_metadata(filters, limit=5000)}]
  Busca datasets CMIP6 disponibles en ESGF mediante filtros. Devuelve \texttt{pd.DataFrame} con metadatos.

  \item[\apisig{get\_all\_urls(dataset\_id)}]
  Obtiene la lista de URLs de descarga para un dataset ESGF. Devuelve lista de cadenas.

  \item[\apisig{get\_combination\_if\_complete(model, experiment, variant, variables)}]
  Verifica que todos los ficheros de una combinación están disponibles en ESGF. Devuelve \texttt{pd.DataFrame} o \texttt{None}.

  \item[\apisig{download\_file(url, local\_path, chunk\_size=1024*1024, max\_retries=3)}]
  Descarga un fichero desde una URL HTTP/OPeNDAP con reintentos. Devuelve \texttt{bool} indicando éxito.

  \item[\apisig{download\_soilgrids(output\_dir, max\_retries=10, include\_metadata=False)}]
  Descarga las capas SoilGrids utilizadas por \hydra (arcilla, arena, limo, materia orgánica). Sin valor de retorno.

  \item[\apisig{open\_soilgrids\_geotif(file)}]
  Lee un GeoTIFF de SoilGrids. Devuelve tupla \texttt{(array, transform, crs)}.

  \item[\apisig{find\_usda\_soilclass(sand, silt, clay, no\_data\_value=255)}]
  Clasifica textura USDA a partir de fracciónes de arena, limo y arcilla. Devuelve \texttt{ndarray} de clases USDA.

  \item[\apisig{create\_new\_tif(template\_file, data, output\_file)}]
  Crea un nuevo GeoTIFF copiando la georreferencia de un fichero plantilla. Sin valor de retorno.

  \item[\apisig{extract\_mode\_soilclass(input\_dir, output\_path)}]
  Calcula la clase de suelo modal entre varias capas de profundidad. Sin valor de retorno.
\end{description}

% ============================================================
\section{\texttt{pyhydra.climate.time\_series}}
% ============================================================

\subsection{Extracción de eventos (\texttt{events})}

\begin{description}[style=nextline,font=\ttfamily]
  \item[\apisig{extract\_events(series, threshold, variable=\textquotesingle{}discharge\textquotesingle{},\\ method=\textquotesingle{}spell\textquotesingle{}, threshold2=None, min\_duration=1, min\_gap=1,\\ min\_sep=7, n\_days=3, plot=False)}]
  Punto de entrada unificado para extracción de eventos. \texttt{variable} acepta \texttt{\textquotesingle{}discharge\textquotesingle{}} o \texttt{\textquotesingle{}precipitation\textquotesingle{}}. \texttt{method} para precipitación: \texttt{\textquotesingle{}spell\textquotesingle{}} (rachas húmedas), \texttt{\textquotesingle{}pot\textquotesingle{}} (picos sobre umbral) o \texttt{\textquotesingle{}nday\textquotesingle{}} (acumulación en $N$ días). Devuelve tupla \texttt{(stats\_df, bounds\_df)}.

  \item[\apisig{extract\_discharge\_events(series, threshold,\\ threshold2=None, plot=False)}]
  Identifica eventos de crecida por cruce de umbral con refinamiento en puntos de inflexión. \texttt{threshold2} es el umbral mínimo de pico para retener un evento. Devuelve tupla \texttt{(stats\_df, bounds\_df)} donde \texttt{stats\_df} tiene columnas \texttt{peak}, \texttt{mean}, \texttt{duration}, \texttt{volume}, \texttt{date\_peak}.

  \item[\apisig{extract\_precipitation\_events(series, threshold=0.0, min\_duration=1,\\ min\_gap=1)}]
  Extrae rachas húmedas continuas. Pasos secos menores que \texttt{min\_gap} se fusionan con los eventos adyacentes. Devuelve tupla \texttt{(stats\_df, bounds\_df)} con columnas \texttt{peak}, \texttt{total}, \texttt{duration}, \texttt{mean\_intensity}, \texttt{date\_start}.

  \item[\apisig{extract\_precipitation\_events\_pot(series, threshold, min\_sep=7)}]
  Extrae eventos extremos por picos sobre umbral con declustering. Expande cada pico a su racha húmeda completa. Devuelve tupla \texttt{(stats\_df, bounds\_df)}.

  \item[\apisig{extract\_precipitation\_events\_nday(series, n\_days=3, threshold=None,\\ min\_sep=None)}]
  Extrae eventos de acumulación en ventana de $N$ días. Útil para cuencas donde la inundación responde a acumulaciones multidia. Devuelve tupla \texttt{(stats\_df, bounds\_df)}.

  \item[\apisig{extract\_concurrent\_events(event\_bounds, series\_dict, buffer\_days=0,\\ stats=(\textquotesingle{}max\textquotesingle{}, \textquotesingle{}mean\textquotesingle{}, \textquotesingle{}total\textquotesingle{}))}]
  Extrae estadísticos en múltiples estaciones durante las ventanas de eventos detectados en una estación objetivo. Útil para análisis de extremos compuestos y coherencia espacial. Devuelve \texttt{pd.DataFrame} con una fila por evento.
\end{description}

\subsection{Análisis de extremos (\texttt{extremes})}

\begin{description}[style=nextline,font=\ttfamily]
  \item[\apisig{extract\_block\_maxima(series, freq=\textquotesingle{}YE\textquotesingle{})}]
  Extrae máximos por bloque temporal. \texttt{freq=\textquotesingle{}YE\textquotesingle{}} para máximos anuales, \texttt{\textquotesingle{}QE\textquotesingle{}} estacionales, \texttt{\textquotesingle{}ME\textquotesingle{}} mensuales. Devuelve \texttt{pd.Series}.

  \item[\apisig{extract\_pot(series, threshold, min\_gap=7)}]
  Extrae picos independientes sobre umbral. \texttt{min\_gap} es el número mínimo de pasos entre picos independientes. Devuelve \texttt{pd.Series} de picos indexada por fecha.

  \item[\apisig{threshold\_stability\_plot(series, thresholds=None, min\_peaks=10)}]
  Diagrama de estabilidad de umbral GPD: exceso medio y parámetro de forma frente al umbral. Devuelve \texttt{pd.DataFrame} con columnas \texttt{threshold}, \texttt{n\_exceed}, \texttt{mean\_excess}, \texttt{gpd\_shape}, \texttt{gpd\_scale}.

  \item[\apisig{fit\_gev(data, method=\textquotesingle{}mle\textquotesingle{}, xi\_bounds=(-0.5, 0.8))}]
  Ajusta una distribución GEV. \texttt{method} acepta \texttt{\textquotesingle{}mle\textquotesingle{}} (robusto con arranque por L-momentos), \texttt{\textquotesingle{}lmom\textquotesingle{}} o \texttt{\textquotesingle{}both\textquotesingle{}}. Devuelve \texttt{dict} con claves \texttt{mu}, \texttt{sigma}, \texttt{xi} (convenio: $\xi > 0$ = cola pesada Fréchet).

  \item[\apisig{fit\_gpd(series, threshold, method=\textquotesingle{}mle\textquotesingle{}, min\_gap=7,\\ xi\_bounds=(-0.5, 1.0))}]
  Ajusta una GPD a los excesos sobre \texttt{threshold}. Devuelve \texttt{dict} con claves \texttt{scale}, \texttt{shape}, \texttt{threshold}, \texttt{n\_exceed}, \texttt{lambda\_rate}.

  \item[\apisig{fit\_gev\_map(data, mu\_prior\_mean=None, sigma\_prior\_scale=0.5,\\ xi\_prior\_std=0.3)}]
  Estimación bayesiana MAP de una GEV. No requiere MCMC; rápido y estable con priores poco informativos. Devuelve \texttt{dict} con claves \texttt{mu}, \texttt{sigma}, \texttt{xi}, \texttt{map\_logpost}.

  \item[\apisig{fit\_gev\_fisher(data, n\_samples=1000)}]
  Incertidumbre paramétrica por aproximación de la información de Fisher (covarianza asintótica). Válido cuando $n \geq 20$ máximos anuales. Devuelve \texttt{pd.DataFrame} con columnas \texttt{mu}, \texttt{sigma}, \texttt{xi}.

  \item[\apisig{fit\_gev\_mcmc(data, n\_samples=2000, n\_chains=4, adapt\_delta=0.95)}]
  Inferencia bayesiana completa GEV mediante MCMC con PyMC y el muestreador NUTS. Parameterización no centrada. Requiere \texttt{pip install pymc}. Devuelve \texttt{pd.DataFrame} con muestras posteriores de \texttt{mu}, \texttt{sigma}, \texttt{xi}.

  \item[\apisig{return\_level\_gev(params, T)}]
  Cuantil GEV para período de retorno $T$ (años). Devuelve \texttt{float}.

  \item[\apisig{return\_level\_gpd(params, T)}]
  Cuantil GPD bajo modelo de proceso de Poisson para período de retorno $T$. Devuelve \texttt{float}.

  \item[\apisig{return\_levels(params, T\_values, dist=\textquotesingle{}gev\textquotesingle{})}]
  Calcula cuantiles para múltiples períodos de retorno. \texttt{dist} acepta \texttt{\textquotesingle{}gev\textquotesingle{}} o \texttt{\textquotesingle{}gpd\textquotesingle{}}. Devuelve \texttt{pd.Series} indexada por período de retorno.

  \item[\apisig{return\_level\_ci(data, T, dist=\textquotesingle{}gev\textquotesingle{}, method=\textquotesingle{}mle\textquotesingle{},\\ n\_bootstrap=500, ci=0.95, threshold=None)}]
  Intervalo de confianza bootstrap para un cuantil de período de retorno. Devuelve tupla \texttt{(estimación, límite\_inferior, límite\_superior)}.

  \item[\apisig{plot\_return\_levels(data, params, dist=\textquotesingle{}gev\textquotesingle{}, T\_range=(1.5, 1000),\\ n\_bootstrap=200, ci=0.95, ax=None)}]
  Gráfico de niveles de retorno con banda de confianza bootstrap. Devuelve \texttt{matplotlib.Axes}.

  \item[\apisig{plot\_diagnostic(data, params, dist=\textquotesingle{}gev\textquotesingle{}, title=None, axes=None)}]
  Panel diagnóstico $2 \times 2$: niveles de retorno, ajuste de densidad, Q-Q y P-P. Devuelve \texttt{matplotlib.Figure}.
\end{description}

\subsection{Métricas de calibración e indicadores de calidad (\texttt{stats})}

\begin{description}[style=nextline,font=\ttfamily]
  \item[\apisig{compute\_nse(obs, sim)}]
  Nash--Sutcliffe Efficiency (NSE) entre series observada y simulada. $\text{NSE} = 1$ indica ajuste perfecto; $\text{NSE} < 0$ indica que la media de la observación es mejor predictor que el modelo. Acepta \texttt{numpy.ndarray} o \texttt{pd.Series} alineadas en tiempo. Devuelve \texttt{float}.

  \item[\apisig{compute\_kge(obs, sim)}]
  Kling--Gupta Efficiency (KGE), métrica compuesta que descompone el error en correlación, sesgo y variabilidad relativa. $\text{KGE} = 1$ indica ajuste perfecto; valores $>0$ son generalmente aceptables. Devuelve \texttt{float}.

  \item[\apisig{compute\_pbias(obs, sim)}]
  Sesgo porcentual (PBIAS) entre observado y simulado. Valores positivos indican subestimación del modelo; negativos, sobreestimación. Devuelve \texttt{float} en porcentaje.
\end{description}

% ============================================================
\section{\texttt{pyhydra.climate.spatial\_analysis}}
% ============================================================

\subsection{Análisis regional de frecuencias (\texttt{rfa})}

\begin{description}[style=nextline,font=\ttfamily]
  \item[\apisig{fit\_gev\_mle(data, xi\_bounds=(-0.5, 0.8))}]
  Ajuste GEV por MML con arranque por L-momentos y parámetro de forma acotado. Devuelve \texttt{dict} con \texttt{mu}, \texttt{sigma}, \texttt{xi}.

  \item[\apisig{fit\_gev\_lmom(data)}]
  Ajuste GEV por el método de L-momentos (requiere \texttt{lmoments3}). Devuelve \texttt{dict} con \texttt{mu}, \texttt{sigma}, \texttt{xi}.

  \item[\apisig{fit\_gev\_bayes(data, n\_chains=4, n\_samples=1000)}]
  Inferencia bayesiana GEV mediante MCMC. Envoltorio de \texttt{fit\_gev\_mcmc} de \texttt{time\_series.extremes}. Devuelve \texttt{pd.DataFrame} con columnas \texttt{mu}, \texttt{sigma}, \texttt{xi}.

  \item[\apisig{return\_level(params, T)}]
  Cuantil GEV para período de retorno $T$ (escalar o array). Devuelve valor(es) de retorno.

  \item[\apisig{return\_level\_bayes(posterior, T, credible=0.95)}]
  Distribución posterior del cuantil de período $T$ a partir de muestras MCMC. Devuelve \texttt{dict} con \texttt{median}, \texttt{lower}, \texttt{upper}.

  \item[\apisig{regional\_index\_flood(data\_dict)}]
  Normaliza cada serie por su índice de avenida (media de máximos anuales). Devuelve tupla \texttt{(normalised\_dict, index\_floods\_series)}.

  \item[\apisig{fit\_regional\_gev(data\_dict, method=\textquotesingle{}lmom\textquotesingle{})}]
  Ajusta una GEV regional sobre datos agrupados normalizados (formulación index-flood). Devuelve tupla \texttt{(regional\_params\_dict, index\_floods\_series)}.

  \item[\apisig{regional\_return\_levels(data\_dict, T\_values=(2,5,10,20,50,100),\\ method=\textquotesingle{}lmom\textquotesingle{})}]
  Calcula cuantiles de período de retorno por estación aplicando escalado por índice de avenida. Devuelve \texttt{pd.DataFrame} (estaciones × períodos de retorno).

  \item[\apisig{HierarchicalGEV(T\_values=(2,10,50,100), n\_chains=4, n\_samples=1000,\\ warmup=1000, adapt\_delta=0.99)}]
  Modelo GEV jerárquico bayesiano para análisis regional. Combina datos de varias estaciones compartiendo parámetros a través de hiperpriors poblacionales. Requiere \texttt{pymc} (o \texttt{pystan} como alternativa). Métodos:

  \begin{itemize}
    \item \texttt{.fit(data\_dict)} — Calibra el modelo (PyMC por defecto, PyStan como fallback). \texttt{data\_dict}: dict de nombre de estación a array de máximos anuales. Devuelve \texttt{self}.
    \item \texttt{.posterior\_summary()} — Resumen de la posterior de parámetros poblacionales y por estación. Devuelve \texttt{pd.DataFrame} con media, desviación típica e IC~95\%.
    \item \texttt{.return\_levels(credible=0.95)} — Cuantiles de período de retorno para cada estación. Devuelve \texttt{pd.DataFrame} con columnas \texttt{T\{T\}\_median}, \texttt{T\{T\}\_lower}, \texttt{T\{T\}\_upper}.
  \end{itemize}
\end{description}

\subsection{Cópulas (\texttt{copulas})}

\begin{description}[style=nextline,font=\ttfamily]
  \item[\apisig{fit\_marginal(data, variable\_name=\textquotesingle{}\textquotesingle{})}]
  Seleccióna la mejor distribución univaríante por BIC (requiere \texttt{openturns}). Devuelve tupla \texttt{(distribución\_ot, muestra\_ot)}.

  \item[\apisig{fit\_discrete\_marginal(data)}]
  Ajusta una distribución empírica discreta a datos de tipo entero. Devuelve tupla \texttt{(distribución\_ot, muestra\_ot)}.

  \item[\apisig{FloodEventCopula(continuous\_vars, discrete\_vars)}]
  Modelización multivariante de eventos de crecida mediante cópula Normal (requiere \texttt{openturns}). Métodos:
  \begin{itemize}
    \item \texttt{.fit(data)} — Ajusta marginales por BIC y cópula Normal. Devuelve \texttt{self}.
    \item \texttt{.sample(n\_samples)} — Genera eventos sintéticos. Devuelve \texttt{pd.DataFrame}.
    \item \texttt{.plot\_marginals()} — Gráficos ECDF vs distribución ajustada.
    \item \texttt{.plot\_dependence(observed, synthetic)} — Scatter plots de pares de variables.
  \end{itemize}

  \item[\apisig{BivariateCopula(family=\textquotesingle{}gumbel\textquotesingle{}, marginal\_families=(\textquotesingle{}gev\textquotesingle{},\textquotesingle{}lognorm\textquotesingle{},\textquotesingle{}gamma\textquotesingle{}))}]
  Análisis de eventos compuestos bivariantes. \texttt{family} acepta \texttt{\textquotesingle{}gumbel\textquotesingle{}}, \texttt{\textquotesingle{}clayton\textquotesingle{}} o \texttt{\textquotesingle{}frank\textquotesingle{}} (no hay selección automática de familia por AIC; sí la hay para cada marginal). Métodos principales:
  \begin{itemize}
    \item \texttt{.fit(x, y, labels=(\textquotesingle{}X\textquotesingle{},\textquotesingle{}Y\textquotesingle{}))} — Ajusta marginales paramétricas (AIC) y el parámetro de la cópula por $\tau$ de Kendall. Devuelve \texttt{self}.
    \item \texttt{.joint\_return\_period(t, kind=\textquotesingle{}or\textquotesingle{})} — Isolínea AND u OR para período $t$ (alias de \texttt{.return\_period\_contour}). Devuelve \texttt{(x\_vals, y\_vals)}.
    \item \texttt{.mpde(t, scenario=\textquotesingle{}OR\textquotesingle{})} — Evento de Diseño de Máxima Probabilidad sobre la isolínea de período $t$ (alias de \texttt{.most\_probable\_event}). Devuelve \texttt{(x\_mpde, y\_mpde)}.
  \end{itemize}

  \item[\apisig{TrivariateCopula(family=\textquotesingle{}gumbel\textquotesingle{}, marginal\_families=(\textquotesingle{}gev\textquotesingle{},\textquotesingle{}lognorm\textquotesingle{},\textquotesingle{}gamma\textquotesingle{}))}]
  Extensión a tres variables del análisis de cópulas Archimedanas exchangeable. \texttt{.fit(x, y, z, labels=(...))} ajusta las tres marginales y un único parámetro de cópula (media de las $\tau$ de Kendall por pares). Misma familia de métodos de período de retorno que \texttt{BivariateCopula}, más \texttt{.conditional\_contours(z\_quantile, T\_list)} y \texttt{.plot(T\_list, z\_quantiles)}.

  \item[\apisig{VineCopula(family\_set=None, marginal\_families=(\textquotesingle{}gev\textquotesingle{},\textquotesingle{}lognorm\textquotesingle{},\textquotesingle{}gamma\textquotesingle{}))}]
  Vine cópulas $d$-dimensionales (R-vine con selección automática de estructura y familia por pareja) para dependencia multivariante flexible (requiere \texttt{pyvinecopulib}). \texttt{.fit(data, labels=None)} ajusta marginales y estructura; \texttt{.sample()}, \texttt{.jcdf()}, \texttt{.kendall\_return\_period()} y \texttt{.most\_probable\_event()} completan la interfaz, siguiendo la formulación de Urrea Méndez et al. (en prensa) \parencite{urrea2026vine}.

  \item[\apisig{GaussianCopulaSampler(data)}]
  Muestreo desde una cópula Gaussiana ajustada a los datos de entrada. Método \texttt{.sample(n)} devuelve \texttt{pd.DataFrame}.
\end{description}

\subsection{Interpolación espacial (\texttt{interpolation})}

\begin{description}[style=nextline,font=\ttfamily]
  \item[\apisig{IDWInterpolator(power=2.0)}]
  Interpolación por distancia inversa ponderada. Sin dependencias externas. Métodos:
  \begin{itemize}
    \item \texttt{.fit(data, x\_col=\textquotesingle{}lon\textquotesingle{}, y\_col=\textquotesingle{}lat\textquotesingle{}, value\_col=\textquotesingle{}value\textquotesingle{})} — Almacena coordenadas y valores. Devuelve \texttt{self}.
    \item \texttt{.predict(grid\_data, x\_col, y\_col)} — Interpola en nuevas localizaciones. Devuelve \texttt{ndarray}.
    \item \texttt{.predict\_timeseries(grid\_data, timeseries\_df, x\_col, y\_col)} — Reconstruye el campo diario para todos los pasos temporales. Devuelve \texttt{pd.DataFrame}.
  \end{itemize}

  \item[\apisig{KrigingInterpolator(variogram\_model=\textquotesingle{}spherical\textquotesingle{}, drift\_terms=None,\\ nlags=6, weight=False)}]
  Kriging universal (requiere \texttt{pykrige}). Métodos:
  \begin{itemize}
    \item \texttt{.fit(data, x\_col, y\_col, value\_col, external\_drift=None)} — Ajusta variograma. Devuelve \texttt{self}.
    \item \texttt{.predict(grid\_data, x\_col, y\_col)} — Devuelve tupla \texttt{(media, varianza)}.
  \end{itemize}

  \item[\apisig{GaussianProcessInterpolator(covariates=None, kernel=None,\\ n\_restarts\_optimizer=5, test\_size=0.2, random\_state=42)}]
  Regresión GP con covariables (requiere \texttt{scikit-learn}). Covariables por defecto: \texttt{['lon', 'lat']}. Métodos:
  \begin{itemize}
    \item \texttt{.fit(data, target\_col)} — Ajusta el GP con validación cruzada. Devuelve \texttt{dict} con \texttt{r2} y \texttt{rmse}.
    \item \texttt{.predict(grid\_data)} — Devuelve tupla \texttt{(media, desviación\_típica)}.
    \item \texttt{.predict\_timeseries(grid\_data, timeseries\_df)} — Ajusta un GP por paso temporal. Devuelve \texttt{pd.DataFrame}.
  \end{itemize}

  \item[\apisig{NeopreneGPReconstructor(gp\_covariates=None,\\ temporal\_resolution=\textquotesingle{}d\textquotesingle{}, seasonality=\textquotesingle{}monthly\textquotesingle{},\\ n\_restarts\_optimizer=3)}]
  Reconstrucción de campo de precipitación diario combinando NEOPRENE y GP. Calibra un modelo NSRP monopuntual por estación y luego interpola los campos sintéticos con GP. Métodos:
  \begin{itemize}
    \item \texttt{.fit(station\_obs, verbose=False)} — Calibra un NSRP por estación.
    \item \texttt{.simulate\_stations(year\_ini, year\_fin)} — Simula series en cada estación.
    \item \texttt{.simulate\_spatial(station\_meta, grid\_meta, year\_ini, year\_fin)} — Genera campo diario espacialmente coherente. Devuelve \texttt{pd.DataFrame} (fechas × celdas de malla).
  \end{itemize}

  \item[\apisig{CopulaCDFEmulator(n\_mc\_samples=1\_000\_000, batch\_size=50\_000)}]
  Emulador GP para la CDF de una cópula multivariante compleja. Permite evaluar la CDF a coste reducido tras un entrenamiento inicial. Métodos:
  \begin{itemize}
    \item \texttt{.fit(copula\_model, n\_dims)} — Entrena el emulador sobre muestras Monte Carlo. Requiere \texttt{openturns}.
    \item \texttt{.predict(query\_points)} — Evalúa la CDF emulada. Devuelve \texttt{ndarray}.
  \end{itemize}
\end{description}

% ============================================================
\section{\texttt{pyhydra.climate.stochastic\_generation}}
% ============================================================

\subsection{Modelos monopuntuales (\texttt{point})}

\begin{description}[style=nextline,font=\ttfamily]
  \item[\apisig{analyze\_ts(series, season=\textquotesingle{}month\textquotesingle{}, dist=\textquotesingle{}gengamma\textquotesingle{},\\ acs\_id=\textquotesingle{}weibull\textquotesingle{}, norm\_type=\textquotesingle{}N1\textquotesingle{}, n\_points=30, lag\_max=30)}]
  Ajusta un modelo estocástico estacional (CoSMoS\_py) a una serie diaria. Para precipitación se recomienda \texttt{NSRPModel}. Devuelve \texttt{dict} del modelo ajustado.

  \item[\apisig{report\_ts(analyzed, method=\textquotesingle{}stat\textquotesingle{})}]
  Resume el modelo ajustado. \texttt{method=\textquotesingle{}stat\textquotesingle{}} devuelve \texttt{pd.DataFrame} con estadísticos estacionales; \texttt{\textquotesingle{}dist\textquotesingle{}} y \texttt{\textquotesingle{}acs\textquotesingle{}} generan gráficos.

  \item[\apisig{simulate\_ts(analyzed, from\_date=None, to\_date=None)}]
  Genera una serie sintética diaria a partir del modelo ajustado. Devuelve \texttt{pd.DataFrame} con columnas \texttt{date} y \texttt{value}.

  \item[\apisig{generate\_ts(n, dist, dist\_params, acs\_values, p0=0.0, n\_series=1)}]
  Generación de bajo nivel a partir de parámetros conocidos. Devuelve lista de \texttt{ndarray}.

  \item[\apisig{fit\_distribution(values, dist=\textquotesingle{}gengamma\textquotesingle{}, norm\_type=\textquotesingle{}N1\textquotesingle{},\\ n\_points=30, opts=None)}]
  Ajusta una distribución marginal a valores positivos. Devuelve \texttt{dict} con \texttt{dist}, \texttt{params\_dict}, \texttt{edf}, \texttt{objective}.

  \item[\apisig{fit\_acs(series, acs\_id=\textquotesingle{}weibull\textquotesingle{}, season=\textquotesingle{}month\textquotesingle{}, lag\_max=30)}]
  Ajusta la estructura de autocorrelación estacional. Devuelve lista de \texttt{dict} de ajuste ACS, uno por temporada.

  \item[\apisig{NSRPModel(temporal\_resolution=\textquotesingle{}d\textquotesingle{}, seasonality=\textquotesingle{}monthly\textquotesingle{},\\ process=\textquotesingle{}normal\textquotesingle{}, statistics=None, weights=None,\\ n\_iterations=100, n\_bees=100, n\_initializations=1,\\ model\_bounds=None, seasonality\_user=None,\\ calibration\_yaml=None)}]
  Generador estocástico monopuntual NSRP (Neyman-Scott de pulsos rectangulares) mediante NEOPRENE. Métodos:
  \begin{itemize}
    \item \texttt{.compute\_statistics(series)} — Calcula estadísticos observados de una \texttt{pd.Series}.
    \item \texttt{.calibrate(series\_or\_stats=None, verbose=False)} — Calibra mediante PSO. Devuelve resultado de calibración NEOPRENE.
    \item \texttt{.fit(series, verbose=False)} — Atajo: calcula estadísticos y calibra en un solo paso.
    \item \texttt{.simulate(year\_ini, year\_fin, simulation\_yaml=None)} — Genera serie sintética. Devuelve resultado de simulación NEOPRENE con atributo \texttt{Daily\_Simulation}.
    \item \texttt{.summary()} — Compara estadísticos observados y ajustados. Devuelve \texttt{pd.DataFrame} o \texttt{None}.
  \end{itemize}
\end{description}

\subsection{Modelos espaciales (\texttt{spatial})}

\begin{description}[style=nextline,font=\ttfamily]
  \item[\apisig{STNSRPModel(temporal\_resolution=\textquotesingle{}d\textquotesingle{}, seasonality=\textquotesingle{}monthly\textquotesingle{},\\ process=\textquotesingle{}normal\textquotesingle{}, statistics=None, weights=None,\\ n\_iterations=100, n\_bees=100, n\_initializations=1,\\ model\_bounds=None, seasonality\_user=None,\\ coordinates=\textquotesingle{}geographical\textquotesingle{}, cell\_radius=None,\\ calibration\_yaml=None)}]
  Generador estocástico multisitio STNSRP. Preserva estadísticos puntuales y correlaciones cruzadas espaciales. El argumento \texttt{statistics} debe incluir \texttt{\textquotesingle{}crosscorr\_1\textquotesingle{}} para el ajuste espacial. \texttt{cell\_radius} son los límites de búsqueda PSO del radio de celda en km (por defecto \texttt{[1.0, 50.0]}). Métodos:
  \begin{itemize}
    \item \texttt{.compute\_statistics(series, attributes)} — Calcula estadísticos y correlaciones.
    \item \path{.calibrate(series_or_stats=None, attributes=None, verbose=False)} — Calibra STNSRP.
    \item \texttt{.fit(series, attributes, verbose=False)} — Calcula estadísticos y calibra en un paso.
    \item \path{.simulate(year_ini, year_fin, attributes, simulation_yaml=None)} — Genera lluvia multisitio espacialmente coherente.
    \item \texttt{.summary()} — Compara estadísticos observados y ajustados.
  \end{itemize}
\end{description}

\subsection{Campos aleatorios (\texttt{fields})}

\begin{description}[style=nextline,font=\ttfamily]
  \item[\apisig{fit\_spatial\_model(spacepoints, p, dist, dist\_params, p0=0.0)}]
  Ajusta un modelo de campo aleatorio espacial (CoSMoS\_py). Devuelve objeto modelo.

  \item[\apisig{generate\_random\_field(n\_steps, model)}]
  Genera $n$ pasos temporales de un campo aleatorio espacial. Devuelve \texttt{ndarray}.

  \item[\apisig{generate\_random\_field\_fast(n\_steps, model)}]
  Versión vectorizada de \texttt{generate\_random\_field} para conjuntos de datos de gran tamaño.

  \item[\apisig{check\_random\_field(field, model)}]
  Verifica que el campo generado reproduce los estadísticos objetivo. Devuelve \texttt{dict} de métricas de validación.

  \item[\apisig{SpatialFieldModel}]
  Clase que encapsula el flujo completo: ajuste, simulación y validación del campo aleatorio espacial.
\end{description}

% ============================================================
\section{\texttt{pyhydra.climate.bias\_correction}}
% ============================================================

\begin{description}[style=nextline,font=\ttfamily]
  \item[\apisig{delta\_method(serie\_raw, serie\_hist, serie\_mod, var, stat)}]
  Corrección de sesgo por método delta. \texttt{var=\textquotesingle{}pr\textquotesingle{}} aplica delta multiplicativo mensual (precipitación); cualquier otro valor aplica delta aditivo (temperatura y otras variables). \texttt{stat} controla la agregación multimodelo: \texttt{\textquotesingle{}mean\textquotesingle{}} o \texttt{\textquotesingle{}median\textquotesingle{}}. Devuelve \texttt{pd.Series} con índice desplazado al período futuro.

  \item[\apisig{BiasCorrection(obs, mod, sce)}]
  Clase con tres métodos de corrección de sesgo:
  \begin{itemize}
    \item \texttt{.quantile\_mapping()} — Mapeado cuantil empírico aditivo.
    \item \texttt{.quantile\_deltamapping()} — Mapeado cuantil empírico multiplicativo (adecuado para precipitación).
    \item \texttt{.scaled\_distribution\_mapping(variable, **kwargs)} — SDM paramétrico: \texttt{variable=\textquotesingle{}precipitation\textquotesingle{}} usa distribución gamma; \texttt{variable=\textquotesingle{}temperature\textquotesingle{}} usa normal.
  \end{itemize}
  Todos los métodos devuelven \texttt{ndarray} con los valores corregidos.
\end{description}

% ============================================================
\section{\texttt{pyhydra.climate.hybrid\_downscaling}}
% ============================================================

\begin{description}[style=nextline,font=\ttfamily]
  \item[\apisig{HydrographClassifier(discharge, events\_bounds, n\_types,\\ n\_points=100)}]
  Clasifica formas de hidrograma mediante remuestreo a longitud fija, PCA y \textit{k}-means. \texttt{n\_types} debe ser un cuadrado perfecto (e.g.\ 9, 16, 25) para disposición en rejilla. Método \texttt{.fit()} (sin argumentos; \texttt{discharge} y \texttt{events\_bounds} ya se pasan en el constructor) devuelve \texttt{pd.DataFrame} con columnas \texttt{Qmax}, \texttt{Qmed}, \texttt{Duracion} y \texttt{shape\_type}.

  \item[\apisig{find\_spatial\_arrangement(n\_rows, n\_cols, centers, iters=200000,\\ method=\textquotesingle{}D8\textquotesingle{})}]
  Organiza los centroides de hidrograma en una rejilla minimizando la distancia entre vecinos. Devuelve rejilla de índices de tipo.

  \item[\apisig{FloodEventSelector(discharge, threshold, threshold2=None,\\ n\_types=25, n\_synthetic=5000, plot=False, output\_dir=None)}]
  Pipeline completo de selección de eventos para downscaling híbrido. Integra extracción de eventos, clasificación de hidrogramas, ajuste de cópula gaussiana y generación de eventos sintéticos. Métodos:
  \begin{itemize}
    \item \texttt{.extract\_events()} — Extrae eventos por cruce de umbral con refinamiento en inflexiones.
    \item \texttt{.classify\_events()} — Clasifica formas de hidrograma.
    \item \texttt{.fit\_copula()} — Ajusta marginales y cópula Gaussiana multivariante.
    \item \texttt{.generate\_synthetic()} — Genera el ensemble sintético.
    \item \texttt{.run(seed=None)} — Ejecuta el pipeline completo en secuencia.
  \end{itemize}

  \item[\apisig{maxdiss(data, n\_select, scalar\_cols, seed\_positions)}]
  Selección de escenarios representativos por máxima disimilitud (MaxDiss). Devuelve tupla \texttt{(submatriz\_seleccionada, posiciones\_elegidas)}.

  \item[\apisig{HydrographReconstructor(centroids, events\_synthetic, discharge\_series,\\ output\_dir=None)}]
  Asigna cada evento sintético a su centroide más próximo y escribe los hidrogramas de entrada para el modelo hidráulico. Método \texttt{.run()} genera ficheros \texttt{Hidrograma\_*.csv}.

  \item[\apisig{FloodMapInterpolator(simulated\_maps, centroids\_sim, k=5,\\ method=\textquotesingle{}knn\textquotesingle{})}]
  Reconstruye mapas de inundación para todos los eventos sintéticos por interpolación k-NN o RBF a partir de las simulaciones de los centroides. Métodos:
  \begin{itemize}
    \item \texttt{.interpolate()} — Reconstruye el ensemble completo de mapas.
    \item \texttt{.compute\_return\_period\_maps(landa, return\_periods)} — Calcula mapas de período de retorno. Devuelve \texttt{dict} de \texttt{T → ndarray}.
  \end{itemize}

  \item[\apisig{FloodMapInterpolatorCC(simulated\_maps, centroids\_sim, k=5,\\ method=\textquotesingle{}knn\textquotesingle{})}]
  Extensión de \texttt{FloodMapInterpolator} para conjuntos con múltiples escenarios climáticos.

  \item[\apisig{pixel\_return\_period(events\_block, landa,\\ return\_periods=(5,10,25,50,100,200,500,1000))}]
  Calcula mapas de período de retorno pixel a pixel. \texttt{events\_block} es un \texttt{ndarray} de forma \texttt{(nrows, ncols, n\_events)} con el ensemble ordenado en el eje 2. \texttt{landa} es la tasa anual de eventos. Devuelve \texttt{dict} mapeando $T$ a \texttt{ndarray(nrows, ncols)}.

  \item[\apisig{save\_return\_period\_geotiffs(calados, template\_tif, output\_dir)}]
  Guarda un GeoTIFF por período de retorno copiando la referencia espacial de la plantilla. Devuelve lista de rutas escritas.
\end{description}

% ============================================================
\section{\texttt{pyhydra.modeling.hydrology}}
% ============================================================

\begin{description}[style=nextline,font=\ttfamily]
  \item[\apisig{generate\_gage(path\_model, name\_gage, file\_gage, names\_gage, ...) }]
  Genera o actualiza el fichero de pluviómetros de HEC-HMS a partir de series temporales. Es una función de escritura de entrada: no ejecuta el modelo y no devuelve datos.

  \item[\apisig{fill\_gage\_series(path\_model, name\_gage, df, variable, ...) }]
  Escribe series de precipitación o caudal en DSS para su uso por HEC-HMS. Se usa como paso previo a \texttt{generate\_met}, \texttt{generate\_run} y la ejecución batch.

  \item[\apisig{generate\_met(path\_model, name\_met, file\_met, names\_gage, ...) }]
  Construye el modelo meteorológico de HEC-HMS enlazando pluviómetros DSS con subcuencas. Sin valor de retorno; modifica ficheros de proyecto.

  \item[\apisig{generate\_hms(path\_model, name\_model, name\_basin, name\_met, ...) }]
  Actualiza la estructura principal del proyecto HEC-HMS y sus referencias a cuenca, meteorología, control y ejecución.

  \item[\apisig{run\_hms\_script(path\_model, name\_model, names\_run, hms\_dir=None)}]
  Ejecuta HEC-HMS por lotes mediante script generado automáticamente. La función encapsula la llamada externa al motor, por lo que la reproducibilidad depende de la versión instalada de HEC-HMS.

  \item[\apisig{read\_hms\_dss\_timeseries(dss\_path, pathname, ...) }]
  Lee series de salida desde ficheros DSS de HEC-HMS y las devuelve como \texttt{pd.Series} o \texttt{pd.DataFrame} según la consulta.

  \item[\apisig{HMSModel(path\_model, name\_model, name\_run, name\_basin,\\ name\_control, time\_interval, pathname, name\_precip,\\ start\_date, end\_date, path\_output, observed=None,\\ parameter\_names=None, hms\_dir=None)}]
  Clase base para ejecutar un proyecto HEC-HMS y devolver el caudal simulado. No calibra por sí sola: actúa como envoltorio de ejecución y lectura de DSS. Métodos principales:
  \begin{itemize}
    \item \texttt{.run\_hms(*params)} — Ejecuta el modelo y devuelve un array de caudal simulado.
    \item \texttt{.\_write\_params(params)} — Punto de extensión para escribir parámetros de calibración.
  \end{itemize}

  \item[\apisig{TextBasinHMSCalibrationModel(..., subbasin, parameters,\\ output\_pathname\_prefix, flow\_unit\_factor=0.028316846592)}]
  Subclase real de \texttt{HMSModel} que calibra parámetros editando texto en el fichero \texttt{.basin}. Se emplea cuando los parámetros HMS pueden modificarse como factores multiplicativos sobre valores base.

  \item[\apisig{calibrate\_hms\_sceua(model, observed, dbname, n\_evals=120,\\ objective=\textquotesingle{}nse\textquotesingle{}, ngs=None)}]
  Calibra un \texttt{TextBasinHMSCalibrationModel} mediante SCE-UA de \texttt{spotpy}. Devuelve el sampler tras la ejecución.

  \item[\apisig{write\_precipitation\_file(df\_coords, df\_series, output\_path,\\ missing\_value=-99.0)}]
  Escribe un fichero de precipitación SWAT clásico \texttt{.pcp} a partir de coordenadas de estación y series diarias.

  \item[\apisig{write\_temperature\_file(df\_coords, df\_tmax, df\_tmin, output\_path,\\ missing\_value=-99.0)}]
  Escribe un fichero de temperatura SWAT clásico \texttt{.tmp} con máximas y mínimas diarias.

  \item[\apisig{edit\_file\_cio(file\_cio\_path, start\_year, end\_year)}]
  Modifica el período de simulación en \texttt{file.cio}. Sin valor de retorno.

  \item[\apisig{write\_swatplus\_precipitation\_files(df\_stations, df\_series,\\ txtinout\_dir, missing\_value=-99.0)}]
  Genera ficheros \texttt{.pcp} individuales y el índice \texttt{pcp.cli} para SWAT+.

  \item[\apisig{write\_swatplus\_temperature\_files(df\_stations, df\_tmax, df\_tmin,\\ txtinout\_dir, missing\_value=-99.0)}]
  Genera ficheros \texttt{.tmp} individuales y el índice \texttt{tmp.cli} para SWAT+.

  \item[\apisig{run\_swat(model\_dir, swat\_exe, timeout=3600)}]
  Ejecuta el binario de SWAT+ en el directorio \texttt{TxtInOut}. Devuelve el código de salida del proceso.
\end{description}

% ============================================================
\section{\texttt{pyhydra.modeling.hydraulic}}
% ============================================================

\begin{description}[style=nextline,font=\ttfamily]
  \item[\apisig{setup\_sfincs\_model(basin\_geom, dem\_path, output\_dir,\\ discharge\_series, src\_points, crs=32630, resolution=100.0,\\ manning=0.04, tref=\textquotesingle{}20200101 000000\textquotesingle{}, ...)}]
  Configura un modelo SFINCS a partir de una geometría de cuenca, un MDE y puntos de entrada de caudal. Devuelve un objeto \texttt{SfincsModel} configurado; el usuario debe llamar a \texttt{mod.write()} para persistir los ficheros.

  \item[\apisig{write\_manning\_wl\_boundary(model\_dir, discharge\_series, ch\_w,\\ ch\_zb, mann\_n=0.035, slope=2e-4, tref=\textquotesingle{}20240101 000000\textquotesingle{})}]
  Añade condición de contorno de nivel de agua y rugosidad al proyecto SFINCS. Escribe ficheros \texttt{sfincs.bnd}, \texttt{sfincs.bzs} y actualiza \texttt{sfincs.inp}.

  \item[\apisig{run\_sfincs(model\_dir, sfincs\_exe, timeout=7200)}]
  Ejecuta el binario SFINCS para el proyecto configurado. Devuelve el código de salida del proceso.

  \item[\apisig{modify\_unsteady\_file(path\_project, name\_project, file\_number,\\ rainfall\_plan\_name, flow\_series, bc\_pathnames)}]
  Crea un nuevo fichero de flujo no permanente HEC-RAS \texttt{.u\#\#} sustituyendo las series de contorno por los datos de \texttt{flow\_series}.

  \item[\apisig{modify\_plan\_file(path\_project, name\_project, plan\_number,\\ rainfall\_plan\_name)}]
  Actualiza un plan HEC-RAS \texttt{.p\#\#} para que apunte al nuevo fichero de flujo no permanente.

  \item[\apisig{modify\_project\_file(path\_project, name\_project, plan\_number,\\ rainfall\_plan\_name)}]
  Cambia el plan activo en el fichero de proyecto \texttt{.prj}.

  \item[\apisig{create\_flow\_series(df, col, window=5)}]
  Suaviza una columna de caudal mediante máximo móvil centrado. Devuelve \texttt{pd.Series}.

  \item[\apisig{run\_hec\_ras(path\_project, name\_project, ras\_version=641)}]
  Ejecuta el plan activo de HEC-RAS mediante \texttt{rascontrol}. Requiere instalación local de HEC-RAS en Windows.

  \item[\apisig{generate\_manning\_combinations(manning\_dist\_csv, n\_samples=1000,\\ mc\_size=10\_000, seed=None)}]
  Genera combinaciones Monte Carlo de coeficientes de Manning para análisis de sensibilidad hidráulica a partir de un CSV de distribuciones bibliográficas por clase de uso del suelo. Devuelve \texttt{pd.DataFrame}.

  \item[\apisig{load\_flood\_ensemble(results\_dir, pattern=\textquotesingle{}hamax\_sim\_*.tif\textquotesingle{},\\ threshold=0.05, chunks=None)}]
  Carga un ensemble de mapas de inundación (GeoTIFF de calado máximo) como \texttt{xarray.DataArray} para análisis estadístico espacial.

  \item[\apisig{load\_sfincs\_ensemble(results\_dir, subdir\_pattern=\textquotesingle{}tmp\_sfincs\_compound\_\{n\}\textquotesingle{},\\ nc\_filename=\textquotesingle{}sfincs\_map.nc\textquotesingle{}, hmax\_var=\textquotesingle{}hmax\textquotesingle{},\\ hmax\_tidx=0, threshold=0.05, crs=32630)}]
  Lee resultados brutos de SFINCS desde subdirectorios de simulación, extrae la variable \texttt{hmax} de cada \texttt{sfincs\_map.nc}, georreferencia las coordenadas \texttt{x}/\texttt{y} y devuelve un \texttt{xarray.DataArray} con dimensión \texttt{simulation}. Esta es la función de lectura de resultados SFINCS disponible actualmente en \pyhydra.

  \item[\apisig{spatial\_stats(ensemble)}]
  Calcula estadísticos espaciales de un ensemble hidráulico, incluyendo área inundada, máximos y métricas resumidas.

  \item[\apisig{manning\_flood\_regression(flood\_ensemble, manning\_ensemble,\\ cell\_area\_m2=25.0, threshold=0.05)}]
  Ajusta modelos de regresión para relacionar parámetros de Manning y métricas de inundación derivadas de los ensembles hidráulicos cargados con \texttt{load\_flood\_ensemble} o \texttt{load\_sfincs\_ensemble}. Devuelve \texttt{pd.DataFrame}.

  \item[\apisig{filter\_anomalous\_simulations(*results, metrics=None,\\ z\_threshold=3.0)}]
  Identifica y filtra simulaciones anómalas comparando una o más tablas de resultados por z-score. Devuelve tupla \texttt{(máscara, resumen)}.
\end{description}

% ============================================================
\section{\texttt{pyhydra.utils}}
% ============================================================

\begin{description}[style=nextline,font=\ttfamily]
  \item[\apisig{interactive\_map(zoom=2, center=(0, 0))}]
  Crea un mapa interactivo Leaflet dentro de JupyterLab mediante \texttt{ipyleaflet}. Útil para visualización exploratoria de geometrías de cuenca, puntos de descarga y mallas de simulación antes de configurar un caso de estudio. Devuelve objeto \texttt{ipyleaflet.Map}.
\end{description}
